
@@M0@@ de @@M1@@~; un
élément @@M2@@ de @@M3@@ détermine un morphisme de foncteurs si, et seulement si, il vérifie :
@@M4@@

Le lecteur se convaincra aisément qu'à des complications de notation près, il suffit de recopier la démonstration de 1.6.

1.8. On notera @@M5@@ la catégorie des @@M6@@-algèbres, et
@@M7@@
le foncteur qui à une @@M8@@-algèbre @@M9@@ associe la @@M10@@-algèbre @@M11@@ munie de la filtration grossière~; on notera
@@M12@@
le foncteur «~oubli de la filtration~». Ces foncteurs vérifient les relations suivantes :
@@M13@@
(cf. 1.3, ex. 2)

\textbf{Lemme 1.9.} \uline{Soit} @@M14@@ \uline{un foncteur sur} @@M15@@~; \uline{alors l'homomorphisme canonique}~:
@@M16@@
